\section{Beyond radial geometry: the companion's scope}\label{art:beyond}

The radial criterion is a useful sufficient condition. It is not necessary
for convergence, and the broader transported certificate is not necessary
either. The complete companion report retains a scalar convergence baseline
and a scalar certificate-nonexistence example that separate these questions.
It also develops three distinct geometric extensions. These additional
results are recorded here for orientation; they are not among the theorems
independently established by this article.

The planar result in the companion's Chapter~9 concerns regular switching
and strict positive-level contours. Its comparison direction moves along
the contour geometry; the proof uses a signed determinant estimate and
one-sided derivative comparison. Its nondegenerate smooth-max application
does not require the radial secant condition. Chapter~11 treats fixed jump
normals in arbitrary dimension, where a smooth-subspace residual budget
replaces the planar argument. Chapter~10 treats a parabolic family with
flat contour pieces; that case does not inherit the strict-contour hypothesis.
The corresponding exact formal contracts and proofs are separate entries in
the companion correspondence index.

These distinctions also limit the geometric interpretation. A certificate
can annihilate accessible limiting gradient jumps under an extension
hypothesis; it need not annihilate every active-gradient difference without
an additional accessibility premise. Affine-criticality describes a geometric
obstruction, while a particular failing run requires its own dynamical
construction. The article uses neither a universal identification of active
and accessible jumps nor an assumed convergence of its metric.

For readers needing quantitative or noncompact conclusions, the companion
also retains the finite-horizon reciprocal bounds, superquadratic reciprocal
sums under summable predicted decreases, noncompact homogeneous alternatives,
affine-fiber lifting, and conservative rates. Their hypotheses and conclusions
are not all ordered by implication. They remain separate from the compact-level
qualitative route used for the article's local positive theorem.
